\section{Discussion}
\label{sec:discussion}

\subsection{Implications for the prevalence literature}
\label{sec:discussion:prevalence}

The measured errors do not resolve into a single multiplier that could be applied to
published figures, and the reason is itself the finding: the error depends on the
stratum. A study that scores long pull request bodies is in a different position from
one that scores commit messages, and the same detector moves from below chance on one
genre to far below chance on another.

The direction of the effect differs by channel. \citet{ji2026exploratory} apply the
zero-shot stack at thresholds imported from other domains. In our corpus those
instruments are below chance in every genre, so their flagged set is not a noisy version
of the AI-written set; it is close to unrelated to it, and the share of artifacts they
flag is closer to a property of the threshold than of the population. The keyword scans
of \citet{mao2026measurement} and \citet{almujahid2026aiusage} fail in the opposite
direction. Their false positive rate in our negative set is zero, so what they count is
real, and their recall on agent-authored text is under four per cent, so what they count
is a small and unrepresentative slice. Their numbers are lower bounds whose distance
from the truth is unknown, which is how they should be read.

Our result and that of \citet{khosravani2026census} bound the problem from opposite
sides. They show that declaration channels undercount, with bot-account lookup
recovering a small fraction of the commits their multi-method approach finds. We show
that content-based detection, offered as the remedy for exactly that gap, does not
recover the difference in this domain. A defensible prevalence figure would therefore
need a lower bound from declaration and an upper bound from corrected detection, and
Table~\ref{tab:youden} says the upper bound is not currently available. Reporting one
number from either channel alone overstates what is known.

\subsection{Implications for platform-level attribution}
\label{sec:discussion:platform}

Declared attribution works where it exists. \AidevTrailerRate{} of commits in AIDev
carry a \texttt{Co-authored-by} trailer, and the agent-account channel is exact by
construction: no inference, no threshold, no error rate to estimate. It is also cheap,
since the platform already knows what opened the pull request.

The contrast with the detection results is the argument for treating that metadata as
infrastructure rather than as a convention. Every artifact whose authorship is declared
at the moment of creation is one that no detector has to guess at later, and the
guessing, as measured here, does not work. The recommendation is narrow and
actionable: preserve agent attribution through merges, squashes and rebases, where it is
currently lost; and prefer declaration over inference wherever both are available. This
does not solve undeclared use, and nothing in this paper suggests it would.

\subsection{Implications for research practice}
\label{sec:discussion:practice}

Table~\ref{tab:checklist} turns the results into requirements for studies that estimate
how much of a corpus is AI-written. Each row names the result behind it, so none of them
rests on our judgement alone.

\begin{table}[t]
\caption{What each result asks of a prevalence study.}
\label{tab:checklist}
\small
\begin{tabular}{@{}p{0.44\textwidth}p{0.46\textwidth}@{}}
\toprule
Requirement & Because \\
\midrule
Validate on the genre being counted before estimating &
Thresholds do not transfer across domains, and every zero-shot cell in
Table~\ref{tab:zeroshot} is below chance \\
\addlinespace
Report per genre, never pooled &
Error rates differ by genre for every instrument (Tables~\ref{tab:zeroshot}
and~\ref{tab:recall}) \\
\addlinespace
Separate pre-language-model automation and report its rate &
An instructed model flags it \JudgeBotRatioMin{}--\JudgeBotRatioMax{} times more often
than human text of the same period (Table~\ref{tab:automation}) \\
\addlinespace
State the operating point with the sensitivity and specificity measured at it &
A share of artifacts flagged is uninterpretable without both \\
\addlinespace
Correct the estimate, and report when the correction is undefined &
Most instrument and genre pairs admit no correction at all
(Table~\ref{tab:estimable}) \\
\bottomrule
\end{tabular}
\end{table}

The last is the one most likely to be resisted, because it converts a headline number
into a statement that no number is available. That is nonetheless what the measurement
supports for most of the instruments in current use.

\subsection{Ethical considerations}
\label{sec:discussion:ethics}

Detector output is already used in governance: contribution policies, review triage and
moderation decisions. The error rates measured here are not evenly distributed. Rates
are highest on short artifacts and on text produced by automation, and
\citet{liang2023biased} report elevated false positives for writing by non-native
English speakers, a population heavily represented among open-source contributors. We
did not measure the linguistic dimension directly, and say so in
Section~\ref{sec:threats}; the design would require repository-level language proxies
and the claim would still be about repositories rather than people.

What the measurement does support is narrower and sufficient. An instrument scoring
below chance in every genre, used as a gate, rejects contributions at a rate unrelated
to whether they were AI-written, and the cost falls on whoever writes text the detector
finds unusual. Deploying these tools for gatekeeping is not supported by their measured
performance in this domain.
